
\subsection{Evaluating \method: Zero-Shot Relation Extraction (zsRE)} \label{subsec:zsre}

We wish to test our localized factual association hypothesis: can storing a single new vector association using \method insert a substantial, generalized factual association into the model?

A natural question is how \method compares to other model-editing methods, which use direct optimization or hypernetworks to incorporate a single new training example into a network. For baselines, we examine Fine-Tuning \textbf{(FT)}, which applies Adam with early stopping at one layer to minimize $-\log \pr{o^* \mid x}$.
Constrained Fine-Tuning \textbf{(FT+L)} \citep{zhu-ft} additionally imposes a parameter-space $L_\infty$ norm constraint on weight changes.
We also test two hypernetworks: \efk \textbf{(KE)} \citep{decao-ke} and \textbf{MEND} \citep{mend}, both of which learn auxiliary models to predict weight changes to $G$. Further details are described in Appendix~\ref{apd:implementations}. 

\begin{wraptable}{R}{0.5\textwidth}
\vspace{-13pt}%
    \centering
    \caption{\small zsRE Editing Results on GPT-2 XL.}
    \vspace{-5pt}%
    \label{tab:zsre-results}    
    \tiny
    \begin{adjustbox}{width=0.5\textwidth}
    \begin{tabular}{l r r r}
        \toprule
        \textbf{Editor} & \textbf{Efficacy} $\uparrow$ & \textbf{Paraphrase} $\uparrow$ & \textbf{Specificity} $\uparrow$ \\
        \midrule
GPT-2 XL & 22.2 ($\pm$0.5) & 21.3 ($\pm$0.5) & 24.2 ($\pm$0.5)\\\midrule
FT & 99.6 ($\pm$0.1) & 82.1 ($\pm$0.6) & 23.2 ($\pm$0.5)\\
FT+L & 92.3 ($\pm$0.4) & \badmetric{47.2 ($\pm$0.7)} & 23.4 ($\pm$0.5)\\
KE & 65.5 ($\pm$0.6) & 61.4 ($\pm$0.6) & 24.9 ($\pm$0.5)\\
KE-zsRE & 92.4 ($\pm$0.3) & 90.0 ($\pm$0.3) & 23.8 ($\pm$0.5)\\
MEND & 75.9 ($\pm$0.5) & 65.3 ($\pm$0.6) & 24.1 ($\pm$0.5)\\
MEND-zsRE & 99.4 ($\pm$0.1) & \goodmetric{99.3 ($\pm$0.1)} & 24.1 ($\pm$0.5)\\
ROME & \goodmetric{99.8 ($\pm$0.0)} & 88.1 ($\pm$0.5) & \goodmetric{24.2 ($\pm$0.5)}\\
        \bottomrule
    \end{tabular}
    \end{adjustbox}
\vspace{-5pt}%
\end{wraptable}
We first evaluate \method on the Zero-Shot Relation Extraction (zsRE) task used in \citet{mend} and \citet{decao-ke}. Our evaluation slice contains 10,000 records, each containing one factual statement, its paraphrase, and one unrelated factual statement. ``Efficacy'' and ``Paraphrase'' measure post-edit accuracy $\smash{\mathbb{I}\big[o^* = \mathrm{argmax}_o \prsub{o}{G^\prime}\big]}$ of the statement and its paraphrase, respectively, while ``Specificity'' measures the edited model's accuracy on an unrelated fact.
Table~\ref{tab:zsre-results} shows the results: \method is competitive with hypernetworks and fine-tuning methods despite its simplicity. We find that it is not hard for \method to insert an association that can be regurgitated by the model.  Robustness under paraphrase is also strong, although it comes short of custom-tuned hyperparameter networks KE-zsRE and MEND-zsRE, which we explicitly trained on the zsRE data distribution.\footnote{Out-of-the-box, they are trained on a WikiText generation task \citep{mend, decao-ke}.}  We find that zsRE's specificity score is not a sensitive measure of model damage, since these prompts are sampled from a large space of possible facts, whereas bleedover is most likely to occur on related \textit{neighboring} subjects.  Appendix \ref{apd:zsre} has additional experimental details.


\subsection{Evaluating \method: Our \cfd Dataset} %
\label{subsec:counterfact-description}

While standard model-editing metrics on zsRE are a reasonable starting point for evaluating \method, they do not provide detailed insights that would allow us to distinguish superficial wording changes from deeper modifications that correspond to a meaningful change about a fact.

In particular, we wish to measure the efficacy of \textit{significant} changes.  \cite{hase2021language} observed that standard model-editing benchmarks underestimate difficulty by often testing only proposals that the model previously scored as likely.  We compile a set of more difficult \textit{false} facts $(s, r, o^*)$: these counterfactuals start with low scores compared to the correct facts $(s, r, o^c)$.
Our Efficacy Score \textbf{(ES)} is the portion of cases for which we have $\mathbb{P}[o^*] > \mathbb{P}[o^c]$  post-edit, and Efficacy Magnitude \textbf{(EM)} is the mean difference $\mathbb{P}[o^*] - \mathbb{P}[o^c]$. Then, to measure \textbf{generalization}, with each counterfactual we gather a set of rephrased prompts equivalent to $(s, r)$ and report Paraphrase Scores \textbf{(PS)} and \textbf{(PM)}, computed similarly to ES and EM. To measure \textbf{specificity}, we collect a set of nearby subjects $s_n$ for which $(s_n, r, o^c)$ holds true. Because we do not wish to alter these subjects, we test $\mathbb{P}[o^c] > \mathbb{P}[o^*]$, reporting the success fraction as Neighborhood Score \textbf{(NS)} and difference as \textbf{(NM)}. To test the generalization--specificity tradeoff, we report the harmonic mean of ES, PS, NS as Score (\textbf{S}).

\begin{wraptable}{R}{0.45\textwidth}
\vspace{-12pt}%
    \centering
    \caption{\small \cfd Composition}
    \label{tab:cfd-summary}    
    \vspace{-6pt}%
    \begin{adjustbox}{width=0.45\textwidth}
    \npdecimalsign{.}
    \nprounddigits{0}
     \begin{tabular}{l r r r}
    \toprule
         & & \multicolumn{1}{c}{\multirow{2}{*}{\shortstack{\bf Per \\ \bf Relation}}} & \multicolumn{1}{c}{\multirow{2}{*}{\shortstack{\bf Per \\ \bf Record}}} \\
         \textbf{Item} & \multicolumn{1}{c}{\textbf{Total}} & & \\ 
        \midrule
        Records & 21919 & 645 & 1 \\
        \midrule
        Subjects & 20391 & 624 & 1 \\
        Objects & 749 & 60 & 1 \\
        Counterfactual Statements & 21595 & 635 & 1 \\
        Paraphrase Prompts & 42876 & 1262 & 2 \\
        Neighborhood Prompts & 82650 & 2441 & 10 \\
        Generation Prompts & 62346 & 1841 & 3 \\
    \bottomrule
    \end{tabular}
    \end{adjustbox}

    \centering
    \caption{\small Comparison to Existing Benchmarks}
    \label{tab:cfd-comparison}
    \vspace{1pt}%
    \begin{adjustbox}{width=0.44\textwidth}
    \begin{tabular}{lccccccc}
    \toprule
        \textbf{Criterion} & SQuAD & zSRE & FEVER & WikiText & \pararel & \textbf{\cfdabbr} \\
        \midrule
        Efficacy & \cmark & \cmark & \cmark & \cmark & \cmark & \cmark \\
        Generalization & \cmark & \cmark & \cmark & \xmark & \cmark & \cmark \\
        Bleedover & \xmark & \xmark & \xmark & \xmark & \xmark & \cmark  \\
        Consistency & \xmark & \xmark & \xmark & \xmark & \xmark & \cmark \\
        Fluency & \xmark & \xmark & \xmark & \xmark & \xmark & \cmark  \\
    \bottomrule
    \end{tabular}
    \end{adjustbox}
\vspace{-12pt}%
\end{wraptable}
We also wish to measure semantic \textbf{consistency} of $\rewg$'s generations. To do so, we generate text starting with $s$ and report \textbf{(RS)} as the $\cos$ similarity between the unigram TF-IDF vectors of generated texts, compared to reference texts about subjects sharing the target property $o^*$.  Finally, we monitor \textbf{fluency} degradations by measuring the weighted average of bi- and tri-gram entropies \citep{Zhang2018GeneratingIA} given by $-\sum_k f(k) \log_2 f(k)$, where $f(\cdot)$ is the $n$-gram frequency distribution, which we report as \textbf{(GE)}; this quantity drops if text generations are repetitive.



In order to facilitate the above measurements, we introduce \cfd, a challenging evaluation dataset for evaluating counterfactual edits in language models. Containing 21,919 records with a diverse set of subjects, relations, and linguistic variations, \cfd's goal is to differentiate robust storage of new facts from the superficial regurgitation of target words. See Appendix \ref{apd:counterfact} for additional technical details about its construction, and Table \ref{tab:cfd-summary} for a summary of its composition.


\subsection{Confirming the Importance of Decisive States Identified by Causal Tracing}
\label{subsec:romeeval-mlp}

\begin{figure}
  \centering
  \includegraphics[keepaspectratio, width=1\textwidth]{figs/rome-sweeps.pdf}
  \vspace{-15pt}%
  \caption{\small \method edits are benchmarked at each layer-and-token combination in GPT-2-XL. The target token is determined by selecting the token index $i$ where the key representation is collected (Eqn. \ref{eq:kstar-sample}). ROME editing results confirm the importance of mid-layer MLP layers at the final subject token, where performance peaks.}
  \vspace{-10pt}%
  \lblfig{rome-sweeps}
\end{figure}


In Section \ref{sec:causal-tracing}, we used Causal Tracing to identify decisive hidden states. To confirm that factual associations are indeed stored in the MLP modules that output those states, we test \method's effectiveness when targeted at various layers and tokens.
\reffig{rome-sweeps} plots four metrics evaluating both generalization (a,b,d) and specificity (c). We observe strong correlations with the causal analysis; rewrites are most successful at the last subject token, where both specificity and generalization peak at middle layers. Targeting earlier \textit{or} later tokens results in poor generalization and/or specificity. Furthermore, the layers at which edits generalize best correspond to the middle layers of the early site identified by Causal Tracing, with generalization peaking at the 18th layer. This evidence suggests that we have an accurate understanding not only of \textit{where} factual associations are stored, but also \textit{how}. Appendix \ref{apd:knowing-saying} furthermore demonstrates that editing the late-layer attention modules leads to regurgitation.

\addtolength{\tabcolsep}{2pt}

\begin{table*}[!t]
    \centering
    \tiny
    \caption{\small \textbf{Quantitative Editing Results}. 95\% confidence intervals are in parentheses. \goodmetric{Green} numbers indicate columnwise maxima, whereas \badmetric{red} numbers indicate a clear failure on either generalization or specificity. The presence of \badmetric{red} in a column might explain excellent results in another. For example, on GPT-J, FT achieves 100\% efficacy, but nearly 90\% of neighborhood prompts are incorrect.}
    \label{tab:cf-results}
    \begin{adjustbox}{width=1\textwidth}
    \begin{tabular}{lrrrrrrrrc}
    \toprule
        \multirow{2.5}{*}{\textbf{Editor}} & \multicolumn{1}{c}{\textbf{Score}} & \multicolumn{2}{c}{\textbf{Efficacy}} & \multicolumn{2}{c}{\textbf{Generalization}} & \multicolumn{2}{c}{\textbf{Specificity}} & \multicolumn{1}{c}{\textbf{Fluency}} & \multicolumn{1}{c}{\textbf{Consistency}} \\
        \cmidrule(lr){2-2}\cmidrule(lr){3-4}\cmidrule(lr){5-6}\cmidrule(lr){7-8}\cmidrule(lr){9-9}\cmidrule(lr){10-10}
        & S $\uparrow$ & ES $\uparrow$ & EM $\uparrow$ & PS $\uparrow$ & PM $\uparrow$ & NS $\uparrow$ & NM $\uparrow$ & GE $\uparrow$ & RS $\uparrow$ \\
        \midrule
GPT-2 XL & 30.5 & 22.2 (0.9) & -4.8 (0.3) & 24.7 (0.8) & -5.0 (0.3) & 78.1 (0.6) & 5.0 (0.2) & 626.6 (0.3) & 31.9 (0.2)\\\midrule
FT & 65.1 & 100.0 (0.0) & 98.8 (0.1) & 87.9 (0.6) & 46.6 (0.8) & \badmetric{40.4 (0.7)} & \badmetric{-6.2 (0.4)} & 607.1 (1.1) & 40.5 (0.3)\\
FT+L & 66.9 & 99.1 (0.2) & 91.5 (0.5) & \badmetric{48.7 (1.0)} & 28.9 (0.8) & 70.3 (0.7) & 3.5 (0.3) & 621.4 (1.0) & 37.4 (0.3)\\
KN & \badmetric{35.6} & \badmetric{28.7 (1.0)} & \badmetric{-3.4 (0.3)} & \badmetric{28.0 (0.9)} & \badmetric{-3.3 (0.2)} & 72.9 (0.7) & 3.7 (0.2) & \badmetric{570.4 (2.3)} & \badmetric{30.3 (0.3)}\\
KE & 52.2 & 84.3 (0.8) & 33.9 (0.9) & 75.4 (0.8) & 14.6 (0.6) & \badmetric{30.9 (0.7)} & \badmetric{-11.0 (0.5)} & \badmetric{586.6 (2.1)} & 31.2 (0.3)\\
KE-CF & \badmetric{18.1} & 99.9 (0.1) & 97.0 (0.2) & 95.8 (0.4) & 59.2 (0.8) & \badmetric{6.9 (0.3)} & \badmetric{-63.2 (0.7)} & \badmetric{383.0 (4.1)} & \badmetric{24.5 (0.4)}\\
MEND & 57.9 & 99.1 (0.2) & 70.9 (0.8) & 65.4 (0.9) & 12.2 (0.6) & \badmetric{37.9 (0.7)} & \badmetric{-11.6 (0.5)} & \goodmetric{624.2 (0.4)} & 34.8 (0.3)\\
MEND-CF & \badmetric{14.9} & \goodmetric{100.0 (0.0)} & \goodmetric{99.2 (0.1)} & \goodmetric{97.0 (0.3)} & \goodmetric{65.6 (0.7)} & \badmetric{5.5 (0.3)} & \badmetric{-69.9 (0.6)} & \badmetric{570.0 (2.1)} & 33.2 (0.3)\\
ROME & \goodmetric{89.2} & 100.0 (0.1) & 97.9 (0.2) & 96.4 (0.3) & 62.7 (0.8) & \goodmetric{75.4 (0.7)} & \goodmetric{4.2 (0.2)} & 621.9 (0.5) & \goodmetric{41.9 (0.3)}\\
\midrule\midrule
GPT-J & 23.6 & 16.3 (1.6) & -7.2 (0.7) & 18.6 (1.5) & -7.4 (0.6) & 83.0 (1.1) & 7.3 (0.5) & 621.8 (0.6) & 29.8 (0.5)\\\midrule
FT & \badmetric{25.5} & \goodmetric{100.0 (0.0)} & \goodmetric{99.9 (0.0)} & 96.6 (0.6) & 71.0 (1.5) & \badmetric{10.3 (0.8)} & \badmetric{-50.7 (1.3)} & \badmetric{387.8 (7.3)} & \badmetric{24.6 (0.8)}\\
FT+L & 68.7 & 99.6 (0.3) & 95.0 (0.6) & \badmetric{47.9 (1.9)} & 30.4 (1.5) & 78.6 (1.2) & \goodmetric{6.8 (0.5)} & \goodmetric{622.8 (0.6)} & 35.5 (0.5)\\
MEND & 63.2 & 97.4 (0.7) & 71.5 (1.6) & \badmetric{53.6 (1.9)} & 11.0 (1.3) & 53.9 (1.4) & \badmetric{-6.0 (0.9)} & 620.5 (0.7) & 32.6 (0.5)\\
ROME & \goodmetric{91.5} & 99.9 (0.1) & 99.4 (0.3) & \goodmetric{99.1 (0.3)} & \goodmetric{74.1 (1.3)} & \goodmetric{78.9 (1.2)} & 5.2 (0.5) & 620.1 (0.9) & \goodmetric{43.0 (0.6)}\\
    \bottomrule
    \end{tabular}
    \end{adjustbox}
\end{table*}%

\addtolength{\tabcolsep}{-2pt}
Table~\ref{tab:cf-results} showcases quantitative results on GPT-2 XL (1.5B) and GPT-J (6B) over 7,500 and 2,000-record test sets in \cfd, respectively.  In this experiment, in addition to the baselines tested above, we compare with a method based on neuron interpretability, Knowledge Neurons \textbf{(KN)}~\citep{dai-kn}, which first selects neurons associated with knowledge via gradient-based attribution, then modifies MLP weights %
at corresponding rows by adding scaled embedding vectors.  We observe that \textbf{all tested methods other than \method exhibit one or both of the following problems}: (F1) overfitting to the counterfactual statement and failing to generalize, or (F2) underfitting and predicting the same new output for unrelated subjects. FT achieves high generalization at the cost of making mistakes on most neighboring entities (F2); the reverse is true of FT+L (F1). KE- and \mbox{MEND-edited} models exhibit issues with both F1+F2; generalization, consistency, and bleedover are poor despite high efficacy, indicating regurgitation. KN is unable to make effective edits (F1+F2). By comparison, \method demonstrates both generalization and specificity.

\subsection{Comparing Generation Results}
\label{subsec:qualitative-analysis}
\begin{figure}[t]%
\centering\includegraphics[width=\columnwidth]{gen-samples/main-paper-samples.pdf}%
\caption{\small \textbf{Comparison of generated text}. Prompts are \textit{italicized}, {\color{ggreen} green} and {\color{red} red} indicate keywords reflecting correct and incorrect behavior, respectively, and {\color{bblue} blue} indicates a factually-incorrect keyword that was already present in $G$ before rewriting. See Section \ref{subsec:qualitative-analysis} for detailed analysis.}
\lblfig{gen-samples}
\vspace{-10pt}
\end{figure}%


\reffig{gen-samples} compares generated text after applying the counterfactual ``\textit{Pierre Curie's area of work is medicine}'' to GPT-2 XL (he is actually a physicist). \textbf{Generalization:} In this case, FT and ROME generalize well to paraphrases, describing the subject as a physician rather than a physicist for various wordings.  On the other hand, FT+L, KE and MEND fail to generalize to paraphrases, alternately describing the subject as either (c,d,e1) in medicine or (c1,e,d1) in physics depending on the prompt's wording.  KE (d) demonstrates a problem with fluency, favoring nonsense repetition of the word \emph{medicine}. \textbf{Specificity:} FT, KE, and MEND have problems with specificity, changing the profession of a totally unrelated subject.  Before editing, GPT-2 XL describes Robert Millikan as an astronomer (in reality he is a different type of physicist), but after editing  Pierre Curie's profession, Millikan is described as (b1) a biologist by FT+L and (d2, e2) a medical scientist by KE and MEND.  In contrast, ROME is specific, leaving Millikan's field unchanged. See 
 Appendix~\ref{apd:gen-samples} for additional examples. 

\subsection{Human evaluation}

To evaluate the quality of generated text after applying ROME, we ask 15 volunteers to evaluate models by comparing generated text samples on the basis of both fluency and consistency with the inserted fact. Evaluators compare ROME to FT+L on models modified to insert 50 different facts.

We find that evaluators are 1.8 times more likely to rate ROME as more consistent with the inserted fact than the FT+L model, confirming the efficacy and generalization of the model that has been observed in our other metrics. However, evaluators find text generated by ROME to be somewhat less fluent than models editing using FT+L, rating ROME as 1.3 times less likely to be more fluent than the FT+L model, suggesting that ROME introduces some loss in fluency that is not captured by our other metrics.  Further details of the human evaluation can be found in Appendix~\ref{apd:human-evaluation}.


\subsection{Limitations}

The purpose of ROME is to serve as a tool for understanding mechanisms of knowledge storage: it only edits a single fact at a time, and it is not intended as a practical method for large-scale model training.  Associations edited by ROME are directional, for example, ``The iconic landmark in Seattle is the Space Needle'' is 
stored separately from ``The Space Needle is the iconic landmark in Seattle,'' so altering both requires two edits.  A scalable approach for multiple simultaneous edits built upon the ideas in ROME is developed in~\citet*{meng2022memit}.

ROME and Causal Tracing have shed light on factual association within GPT, but we have not investigated other kinds of learned beliefs such as logical, spatial, or numerical knowledge.  Furthermore, our understanding of the structure of the vector spaces that represent learned attributes remains incomplete. Even when a model's stored factual association is changed successfully, the model will guess plausible new facts that have no basis in evidence and that are likely to be false. This may limit the usefulness of a language model as a source of facts.



